\documentclass[12pt]{article}
\usepackage[utf8]{inputenc}
\usepackage[T1]{fontenc}
\usepackage{hyperref}
\title{Ambel Language Resources for K-12 Curriculum Development}
\author{Compiled by Claude Sonnet 5.5 (claude-sonnet-5-5) for Andr\'as Kornai}
\date{}
\begin{document}
\maketitle

\section{Introduction}
Ambel (ISO 639-3 wgo; Glottocode waig1244) is an Austronesian language of the Raja Ampat Islands, spoken on Waigeo in West Papua, Indonesia. It belongs to the South Halmahera--West New Guinea branch. One source in the search puts the speaker count at about 1,600, while the survey figure is about 300. The two sources do not agree. Speakers are shifting to Papuan Malay, and few people born after 2000 know the language. There are two dialects, Metsam Ambel and Metnyo Ambel.

\section{Primary Educational Resources}
No primers, textbooks or school materials were found in the search results.

\section{Linguistic and Reference Materials}
Laura Arnold's 2018 University of Edinburgh PhD thesis, \emph{A Grammar of Ambel}, is the main reference and the first major documentation of the language. See \url{https://era.ed.ac.uk/handle/1842/31120} and \url{https://era.ed.ac.uk/items/d26ecd9e-b540-478f-a90b-242300ec8558/full}. A record is also at \url{https://researchportalplus.anu.edu.au/en/publications/a-grammar-of-ambel-an-austronesian-language-of-raja-ampat-west-ne/}.

The Journal of the International Phonetic Association has an Ambel illustration by the same author: \url{https://www.cambridge.org/core/product/D9E2B1D9B3479E646D53D60112712FFD/core-reader}.

Other overviews: \url{https://en.wikipedia.org/wiki/Ambel_language} and \url{https://www.omniglot.com/writing/ambel.htm}.

\section{Religious and Community Resources}
No scripture or community organisation resources were found.

\section{Digital and Audio Resources}
Arnold's documentation data are reported to be archived with the Endangered Languages Archive and with the Center for Endangered Languages Documentation at Universitas Papua. No direct archive URL was seen in the search results.

\section{Curriculum Development Context}
Ambel is not known to be taught in schools, and no pedagogical materials were found. The Arnold grammar and the archived recordings are the starting point. A curriculum would need teaching materials written with Ambel speakers on Waigeo, and with Universitas Papua. The language's tone system should be taken into account when an orthography is chosen.
\end{document}